\section{مقدمه}
مسئله بسته‌بندی در سطل‌ها (\lr{Bin Packing Problem} یا به اختصار \lr{BPP}) مسئله‌ای کلاسیک در بهینه‌سازی ترکیبیاتی است: مجموعه‌ای از اقلام با وزن‌های مثبت داده می‌شود و باید همه آن‌ها را، بدون تجاوز از ظرفیت ثابت هر سطل، در کمترین تعداد سطل قرار داد.

در این گزارش گونه‌ای از مسئله بررسی می‌شود که علاوه بر قید ظرفیت، قیود هم‌مکانی سخت یا \lr{must-link} دارد. هر رابطه چسبندگی میان دو قلم الزام می‌کند که آن دو در یک سطل قرار گیرند. چون هم‌مکانی یک رابطه تراگذری است، تمام اقلام واقع در یک مؤلفه همبند گراف چسبندگی ناگزیر در یک سطل مشترک قرار می‌گیرند. این تعریف باید از ترجیح نرم برای نزدیک‌بودن اقلام، و نیز از قیود ناسازگاری که هم‌مکانی را منع می‌کنند، تفکیک شود.

زمان‌بندی گردش‌کارهای کانتینری و تخصیص منابع در کوبرنتیز نمونه‌ای از زمینه‌هایی است که در آن نگاشت وظایف به منابع و سازوکارهای اشتراک داده مطرح‌اند \cite{C1_Shan2022}؛ همچنین ادبیات زمان‌بندی کانتینرها اهداف گوناگونی مانند بهره‌وری منابع، زمان پاسخ و ملاحظات کربنی را بررسی می‌کند \cite{C4_Yang2025}. با این حال، این منابع مدل یک‌بعدیِ دقیق این گزارش یا الزام هم‌مکانی ناشی از ترافیک شبکه را اثبات نمی‌کنند. مدل حاضر تنها در ساده‌سازی‌ای قابل استفاده است که منابع همگن و تک‌ظرفیتی باشند، گروه‌های هم‌مکانی از پیش و به‌صورت سخت تعیین شده باشند، و هدف کمینه‌کردن تعداد منابع فعال باشد.

از دید محاسباتی، contraction هر مؤلفه چسبندگی به یک قلم با وزن مجموع، نمونه را دقیقاً به یک \lr{BPP} کلاسیک تبدیل می‌کند. بنابراین این گونه در بدترین حالت قویاً \lr{NP-hard} است، زیرا حالت $E=\emptyset$ همان \lr{BPP} کلاسیک است \cite{GareyJohnson1979}؛ اما در مدل تعریف‌شده الزاماً از \lr{BPP} سخت‌تر نیست، چون پس از یک پیش‌پردازش چندجمله‌ای همان مسئله روی اقلام تجمیع‌شده حل می‌شود. در ادامه، این هم‌ارزی مبنای مدل‌سازی و طراحی الگوریتم قرار می‌گیرد؛ کیفیت جواب‌های الگوریتمی نیز باید با کران یا جواب بهینه سنجیده شود و صرف اجرای محاسباتی، اثبات بهینگی محسوب نمی‌شود.
